\section{Library Architecture}
\label{sec:architecture}

The architecture of the library is organised around a clear separation between the logical data model, the intermediate representation of queries, and the backend-specific execution layer. That separation is central to the goal of the project: changing the database should require minimal change in application code and, ideally, only a change in the connector type. To achieve this, the architecture is built from several layers, each with its own responsibility and with strictly delimited interfaces.

The theoretical chapter showed that unification is possible at the level of the logical model and the query IR, but not by ignoring the differences between backend families. The present chapter takes the next step: it shows how that theoretical position is turned into a concrete architecture. The library is therefore treated not merely as an idea, but as an organised system of types, contracts, and execution paths.

\subsection{Architectural goals and constraints}

The architecture was designed under several constraints acting simultaneously. First, the user-facing API must remain within idiomatic C++, without introducing a separate DSL and without relying on runtime reflection. Second, the query structure must be representable as a \code{constexpr} object so that a significant portion of validation can be performed by the compiler. Third, the backend layer must remain extensible without inheritance from a common virtual interface, because such an interface would hide important backend differences and move correctness checks toward runtime.

A fourth constraint follows from the multi-backend nature of the system itself. The library must support SQL and non-SQL connectors without false equivalence. This means that the architecture must simultaneously provide a shared API and express the limitations of a concrete backend in a formal and testable manner. The choice of a trait-based connector layer together with a capability mechanism follows naturally from this requirement.

\begin{figure}[H]
\centering
\begin{tikzpicture}[node distance=1.3cm and 1.4cm]
    \node[ormaccent] (app) {\textbf{Application layer}\\entity types and query composition};
    \node[ormblock, below=of app] (entity) {\textbf{Entity layer}\\\code{property}, \code{relationship}, \code{table\_name\_trait}};
    \node[ormblock, below=of entity] (ir) {\textbf{Query-IR layer}\\\code{field}, \code{Rule}, \code{select\_query}, \code{insert\_query}};
    \node[ormblock, below=of ir] (connector) {\textbf{Connector layer}\\\code{connector\_trait<DB>} and capability tags};
    \node[ormblock, below=of connector] (driver) {\textbf{Driver / runtime layer}\\SQL/BSON/Redis/Cypher/CQL materialisation};

    \draw[ormline] (app) -- (entity);
    \draw[ormline] (entity) -- (ir);
    \draw[ormline] (ir) -- (connector);
    \draw[ormline] (connector) -- (driver);
\end{tikzpicture}
\caption{Layered architecture of the library}
\label{fig:architecture_layers}
\end{figure}

Figure~\ref{fig:architecture_layers} shows that the library follows a consistent vertical organisation. Application code does not communicate directly with the driver; it works with C++ types and query builders. The driver, in turn, does not know domain semantics, but only the materialised form generated by the connector layer. This is precisely why the architecture can remain both strict and extensible.

\subsection{Layered organisation and responsibility separation}

The topmost layer is the application code, which defines entity types, constructs queries through \code{select}, \code{insert}, \code{update}, and \code{deleteq}, and submits them to \code{db<DB>}. The next layer is the logical model of the library --- the entity descriptions and the query IR. The third layer is the connector abstraction, where \code{connector\_trait<DB>} materialises the same logic into backend-specific syntax and behaviour. At the bottom stands the concrete driver, wire protocol, or mock implementation.

This architecture has an important practical effect. If the user replaces \code{SQLiteDB} with \code{MongoDB}, the query IR remains logically the same, but the connector interprets it as filter/projection execution rather than SQL rendering. This does not mean that every IR is admissible on every backend; it means that the \emph{form of abstraction} is shared, while admissibility is determined formally through the capability mechanism.

\begin{table}[H]
\centering
\small
\caption{Principal architectural layers and their responsibilities}
\label{tab:architecture_responsibilities}
\begin{tabularx}{\textwidth}{L{2.8cm}L{4.0cm}Y}
\toprule
\textbf{Layer} & \textbf{Main artefacts} & \textbf{Responsibility} \\
\midrule
Application & entity types, query-builder calls & describes domain logic and formulates operations in idiomatic C++ \\
Logical & \code{property}, \code{relationship}, \code{field}, \code{Rule} & models data and query structure independently of a concrete backend \\
Connector & \code{connector\_trait<DB>}, capability tags & checks admissibility and materialises the IR into backend behaviour \\
Runtime/driver & concrete connection handles and driver APIs & executes the materialised query and returns results \\
\bottomrule
\end{tabularx}
\end{table}

Table~\ref{tab:architecture_responsibilities} shows that the library does not rely on a single monolithic object that does everything. Instead, decisions are localised in separate layers. This matters both for extensibility and for academic analysis, because each claim can be traced to concrete types and modules.

\subsection{The entity layer as architectural foundation}

The entity layer is built around \code{property<T, Name>}, \code{relationship<...>}, and \code{table\_name\_trait<T>}. It plays a dual role. On the one hand, it represents the C++ domain model. On the other hand, it serves as a logical description of the persistent schema or structure. This means that the architecture does not construct a separate runtime metadata registry. Instead, the type system contains enough information to derive field names, value types, relationships, and the target storage container.

Especially important is the fact that the entity layer is independent of the database family. Neither \code{property} nor \code{relationship} is SQL-specific. They describe logical attributes and logical relations. This creates the foundation for comparable behaviour across relational and non-relational backends. At the architectural level, this is exactly what makes it possible for the thesis to argue that the C++ model is the primary carrier of the logical schema.

\subsection{Query IR as the central architectural component}

The query IR is the central architectural component. It combines field references, predicates, placeholders, and additional clauses such as \code{order\_by}, \code{group\_by}, and \code{limit}. Instead of working with textual queries, the library constructs typed objects whose template parameters encode the logical structure of the operation.

From the architectural point of view, this has several consequences. First, structural validation occurs early. Second, the query IR can be analysed by the connector layer without losing semantic information. Third, prepared queries can preserve not an SQL string, but the intermediate representation itself, which is more natural for a library oriented around the C++ type system.

Fourth, the query IR is the single zone in which entity description and later backend materialisation meet. In that sense it is the architectural hinge of the whole library: high enough to remain storage-agnostic, yet concrete enough to be executed.

\subsection{The connector layer and the formal contract}

Formally, the connector layer is implemented through specialisation of \code{connector\_trait<DB>}. This approach is analogous to standard trait templates such as \code{std::iterator\_traits}. The type \code{DB} may be a simple wrapper over a connection handle, while the entire ORM logic remains in the trait specialisation. In this way, the library does not require the user or the connector author to inherit from a common base class.

\code{connector\_trait<DB>} defines at least three central aspects: \code{wire\_type}, \code{cursor\_type}, and \code{execute(...)}. The specialisation may also declare capability tags such as \code{supports\_joins}, \code{supports\_transactions}, \code{supports\_aggregation}, or \code{supports\_concurrent\_execute}. The presence or absence of these pseudo-markers is sufficient for the upper layers to impose compile-time restrictions.

\begin{lstlisting}[language=C++,caption={Excerpt from the formal connector contract},label={lst:connector_contract}]
template <typename DB>
concept is_connector = requires {
    typename connector_trait<DB>::template wire_type<int>::type;
    typename connector_trait<DB>::cursor_type;
};

template <typename DB, typename Cap>
inline constexpr bool has_capability =
    detail::capability_check<DB, Cap>::value;
\end{lstlisting}

Listing~\ref{lst:connector_contract} shows that the contract is structural rather than inheritance-based. This is a decisive architectural choice. If a virtual interface existed with runtime dispatch, some of the type-checkable information would be lost. The library instead prefers a compile-time structure in which the requirements remain visible both to the compiler and to the connector author.

\begin{figure}[H]
\centering
\begin{tikzpicture}[node distance=1.35cm and 1.25cm]
    \node[ormaccent] (query) {\textbf{Typed query}\\join, transaction, aggregation intent};
    \node[ormblock, below=of query] (db) {\code{db<DB>} facade};
    \node[ormblock, below left=of db] (capok) {Capability present\\query remains admissible};
    \node[ormdanger, below right=of db] (capfail) {Capability missing\\compile-time rejection};
    \node[ormblock, below=of capok] (exec) {\code{connector\_trait<DB>::execute}};

    \draw[ormline] (query) -- (db);
    \draw[ormline] (db) -- (capok);
    \draw[ormline] (db) -- (capfail);
    \draw[ormline] (capok) -- (exec);
\end{tikzpicture}
\caption{Capability gating between the query IR and the connector layer}
\label{fig:capability_gating}
\end{figure}

Figure~\ref{fig:capability_gating} is important because it shows where the architecture decides admissibility. This is not the driver and not a runtime log. The decision is taken in the public API layer, which already has access both to the query IR and to the trait-level information about the connector.

\begin{table}[H]
\centering
\small
\caption{Role of the principal capability tags in the architecture}
\label{tab:capability_tags}
\begin{tabularx}{\textwidth}{L{3.6cm}Y}
\toprule
\textbf{Capability} & \textbf{Architectural meaning} \\
\midrule
\code{supports\_joins} & admits join-based query forms for the selected backend \\
\code{supports\_transactions} & allows use of \code{transaction\_guard} together with \code{begin}/\code{commit}/\code{rollback} hooks \\
\code{supports\_aggregation} & marks support for aggregation operations in the logical query layer \\
\code{supports\_embedding} & makes native support for embedded structures explicit \\
\code{supports\_upsert} & marks backends for which upsert belongs to the natural execution contract \\
\code{supports\_bulk\_insert} & denotes the existence of a more efficient bulk-write path \\
\bottomrule
\end{tabularx}
\end{table}

\subsection{The user-facing \code{db<DB>} facade}

The class \code{db<DB>} is the public facade layer. It encapsulates a reference to a concrete backend object and provides a unified execution API. In practice, this layer has two roles. The first is convenience for the user: the library is used through the same interface regardless of backend. The second is architectural filtering: this is precisely where compile-time checks for capabilities and backend compatibility are performed.

Consequently, \code{db<DB>} is not merely a thin wrapper over \code{execute}. It is the place where the library formalises the boundary between admissible and inadmissible behaviour for the selected backend. It is exactly this layer that enables error messages such as ``this connector does not support \code{JOIN} operations'' instead of allowing such errors to surface later as ambiguous runtime behaviour.

\begin{lstlisting}[language=C++,caption={Excerpt from the \code{db<DB>} facade layer},label={lst:db_facade}]
template <typename DB>
    requires is_connector<DB>
class db
{
public:
    template <typename Query>
    auto operator<<(Query q)
    {
        if constexpr (is_select_query<Query>)
        {
            if constexpr (decltype(q.join_clauses())::size > 0)
            {
                static_assert(has_capability<DB, cap::supports_joins>,
                    "orm::db: this connector does not support JOIN operations.");
            }
        }
        return connector_trait<DB>::execute(*conn_, std::move(q));
    }
};
\end{lstlisting}

Listing~\ref{lst:db_facade} shows why \code{db<DB>} is more than a convenient facade. It is the formal location where query structure and connector capability information meet. This is where the compile-time architecture becomes a real limitation on API use.

\subsection{Migrations as an architectural extension}

The architecture also includes a prototype migration layer through \code{migrate<DB>}, \code{live\_schema}, \code{entity\_meta}, and DDL operations such as \code{create\_table\_op}, \code{add\_column\_op}, \code{drop\_column\_op}, and \code{alter\_column\_type\_op}. This layer is particularly important from the standpoint of the thesis because it makes explicit the relation between the C++ model and the real storage schema.

The migration layer is not uniformly complete for every backend, but architecturally it demonstrates that the library is not limited to query execution. With a sufficiently rich \code{connector\_trait<DB>} specialisation, entity description can also participate in schema evolution. This changes the very meaning of the ORM layer: it becomes not only an interface to data, but a potential source of schema-aware tooling.

\subsection{Architectural invariants}

On the basis of the analysed layers, several invariants can be formulated. First, application code should never depend directly on driver APIs. Second, query logic is formulated in a typed IR rather than in backend-native strings. Third, extension toward a new backend proceeds through trait specialisation and capability declarations rather than through changes to the query API itself. Fourth, if an operation is inadmissible for a backend, this should become visible as early as possible --- ideally at compile time.

These invariants make the architecture both engineering-sound and academically defensible. They allow every future extension to be judged not only by whether it ``works'', but by whether it preserves the already formulated layers and boundaries.

\subsection{Architectural conclusion}

In summary, the architecture of the library separates the problem into three levels: logical model, typed query IR, and backend materialisation. This separation enables strict type discipline, formal extensibility, and controlled behaviour across distinct storage models. The next chapter examines this architectural decision more concretely through the implementation of the core subsystems, where it becomes clear how the architectural plan is turned into executable C++ code.
